\providecommand{\iquest}[1]{¿#1}
\section{Gu\'ia de observaci\'on y taller}

\textbf{C\'odigo del instrumento:} INS-OBS-01. \textbf{Dirigido a:} MES-01, COC-01 y al grupo GO-01.
\textbf{Duración estimada:} 45 minutos de observaci\'on y 60 minutos de taller. Se aplica en el local con autorizaci\'on previa del responsable.

\subsection*{Ficha de sesi\'on}
Fecha: \rule{2.8cm}{0.4pt}\quad Turno: \rule{2.2cm}{0.4pt}\quad \'Area observada: \rule{3cm}{0.4pt}\quad
Observador: \rule{3cm}{0.4pt}\quad C\'odigos de roles presentes: \rule{4cm}{0.4pt}.

\subsection*{Lista de observaci\'on}
Para cada \'item se registra S\'i, No o No observado y una nota breve sin nombres propios.

\begin{enumerate}[leftmargin=18pt,itemsep=1pt]
    \item El pedido se dice en voz alta y se repite para confirmación.
    \item Existe un soporte f\'isico o digital donde se anota el pedido.
    \item La cocina recibe los pedidos en el mismo orden en que llegaron.
    \item Se corrige alg\'un pedido por falta de detalle o por duplicidad.
    \item Se consulta el inventario antes de usar un insumo cr\'itico.
    \item Se registra la salida de insumos en el momento del uso.
    \item El cajero recibe el detalle de lo consumido antes de cerrar.
    \item Se genera un resumen de ventas sin transcripci\'on adicional.
    \item Una sugerencia autom\'atica, si existiera, requerir\'ia confirmación humana.
    \item Se identifica al menos un momento de espera atribuible a la comunicaci\'on.
\end{enumerate}

\subsection*{Din\'amica del taller}
\begin{enumerate}[leftmargin=18pt,itemsep=0pt]
    \item \textbf{Contexto (10 min):} el responsable describe el objetivo del taller y se confirma el uso an\'onimo de lo registrado.
    \item \textbf{Recorrido del pedido (25 min):} en una hoja por turnos se dibuja cada paso, qui\'en lo ejecuta y qu\'e informaci\'on entrega.
    \item \textbf{Puntos de control (15 min):} se marcan con una cruz los pasos donde se detectan errores y con una estrella los que requieren decisi\'on humana.
    \item \textbf{Priorizaci\'on (10 min):} cada participante dispone de tres votos para señalar los puntos cr\'iticos.
\end{enumerate}

\subsection*{Cierre y custodia}
El facilitador lee en voz alta los acuerdos, los códigos de rol y la fecha. El original de las notas se guarda en la carpeta institucional privada \texttt{@uteq.edu.ec}; en el repositorio p\'ublico solo se registra su huella SHA-256 en \texttt{evidencias/MANIFIESTO.md}.
